\documentclass[10pt]{article}
\usepackage[letterpaper,margin=0.76in]{geometry}
\usepackage{amsmath,amssymb,amsthm,booktabs,microtype}
\usepackage[T1]{fontenc}
\usepackage{lmodern}
\usepackage[hidelinks]{hyperref}

\newtheorem{theorem}{Theorem}
\newtheorem{lemma}[theorem]{Lemma}
\newtheorem{remark}[theorem]{Remark}
\newcommand{\Z}{\mathbb Z}
\newcommand{\one}{\mathbf 1}
\newcommand{\ip}[2]{\langle #1,#2\rangle}
\newcommand{\zm}{z_{\rm mid}}
\newcommand{\Ax}{A_x}
\newcommand{\Qx}{\mathcal Q_x}
\setlength{\emergencystretch}{3em}

\title{Exact Adjoint Diagonal Return and Hard-Boundary Decomposition\\
for the Literal V59 Haar Lane}
\author{Liang Wang\\
\small School of Electronic Information and Communications, Huazhong University of Science and Technology (HUST),\\
\small Wuhan, China}
\date{August 26, 2026}

\begin{document}
\maketitle

\begin{abstract}
We push a coefficient-independent ordered-rank Haar vector through the adjoint
of the literal V59 prime-shell operator.  A complete-lattice unit-centered row
vanishes exactly by the previously frozen band-limited Poisson theorem once
$H>2Q$.  This does not annihilate the physical row: deleting its diagonal
returns a $B_Q$-weighted midpoint moment of the literal coefficient $\beta$,
while truncation to $(x/2,x]\cap\mathbb Z$ and the jump between the two rank
children create two explicit boundary lanes.  We retain the input and output
unit masks and show that the output correction is jointly centered only.  The
result is an exact normal form for $\ip{\zm}{\Ax\beta}$, not an estimate.  In
particular, no sign, nonzero value, power saving, $L^2$ theorem, or twin-prime
conclusion is asserted.  Exact Gaussian-rational fixtures and 192 stress
families reproduce the finite algebra but are not asymptotic evidence.
\end{abstract}

\section{Frozen operator and ordered-rank test}

Retain the literal V59 scales and coefficient
\begin{equation}
 I_x=(x/2,x]\cap\Z,\quad H=x^{21/32},\quad Q=x^{1/3},\quad
 \Qx=\{q\text{ prime}:Q<q\leq2Q\},                    \label{eq:scales}
\end{equation}
\begin{equation}
 K_H(h)=\widehat\psi_+(h/H),\qquad
 \beta(t)=\frac{\Lambda(t)}{\log t}
 -\sum_{\substack{d\mid t\\d^{400}\leq x^{133}}}\mu(d),       \label{eq:beta}
\end{equation}
where $\operatorname{supp}\psi_+\subset[-1,1]$ and
$K_H(0)=\int\psi_+=1$.  For $u,t\in I_x$, the source-frozen operator is
\begin{equation}
 \Ax(u,t)=\one_{u\ne t}\sum_{q\in\Qx}q\one_{q\nmid u}\one_{q\nmid t}
 K_H(u-t)\left(\one_{u\equiv t\,(q)}-\frac1{q-1}\right).        \label{eq:operator}
\end{equation}
This orientation and the associated scalar are fixed in TPC-247
\cite{tpc247}; the physical V59 source is \cite{v59}.

Write $I_x=\{n_1<\cdots<n_N\}$, assume $N\geq2$, and put
\begin{equation}
 \ell=\lfloor N/2\rfloor,\quad r=N-\ell,
 \quad L=\{n_1,\ldots,n_\ell\},\quad R=I_x\setminus L,           \label{eq:rank}
\end{equation}
\begin{equation}
 \rho^2=\frac{\ell r}{N},\qquad
 \zm=\rho\left(\frac{\one_L}{\ell}-\frac{\one_R}{r}\right).  \label{eq:haar}
\end{equation}
The split is by ordered rank for every real $x$.  The integer-only
$\lfloor3x/4\rfloor$ crosswalk is not used at nonintegral clocks
\cite{tpc253}.

\section{Complete row and exact adjoint decomposition}

For $q\in\Qx$ and $q\nmid t$, define the complete periodic unit row
\begin{equation}
 v_{q,t}(u)=\one_{q\nmid u}
 \left(\one_{u\equiv t\,(q)}-\frac1{q-1}\right),\qquad u\in\Z. \label{eq:v}
\end{equation}
Its mean over a period is zero.  Introduce the adjoint-oriented quantities
\begin{align}
 P^*_{q,t}&=\sum_{u\in\Z}\overline{K_H(u-t)}v_{q,t}(u),          \label{eq:p}\\
 E^*_{q,t}&=\sum_{u\notin I_x}\overline{K_H(u-t)}v_{q,t}(u),    \label{eq:e}\\
 J^*_{q,t}&=\sum_{u\in I_x}\overline{K_H(u-t)}v_{q,t}(u)
                   (\zm(u)-\zm(t)).                             \label{eq:j}
\end{align}
Absolute convergence follows from the Schwartz decay of $K_H$.

\begin{theorem}[Adjoint diagonal--boundary compiler]\label{thm:compiler}
For every real $x$ with $N\geq2$,
\begin{align}
 (\Ax^*\zm)(t)=\sum_{q\in\Qx}q\one_{q\nmid t}
 \bigg[&\zm(t)P^*_{q,t}-\zm(t)E^*_{q,t}+J^*_{q,t} \notag\\
 &-\frac{q-2}{q-1}\overline{K_H(0)}\zm(t)\bigg].               \label{eq:full}
\end{align}
If $H>2Q$, the source-backed complete-lattice theorem gives
$P^*_{q,t}=0$, and hence
\begin{equation}
 (\Ax^*\zm)(t)=\sum_{q\in\Qx}q\one_{q\nmid t}
 \left[-\zm(t)E^*_{q,t}+J^*_{q,t}-\frac{q-2}{q-1}\zm(t)\right].\label{eq:reduced}
\end{equation}
Moreover,
\begin{equation}
 J^*_{q,t}=\begin{cases}
 -\rho^{-1}\displaystyle\sum_{u\in R}\overline{K_H(u-t)}v_{q,t}(u),&t\in L,\\[2mm]
 +\rho^{-1}\displaystyle\sum_{u\in L}\overline{K_H(u-t)}v_{q,t}(u),&t\in R.
 \end{cases}                                                     \label{eq:jump}
\end{equation}
\end{theorem}

\begin{proof}
Fix $q\nmid t$ and abbreviate
$F(u)=\overline{K_H(u-t)}v_{q,t}(u)$.  Adding and subtracting $\zm(t)$ gives
\[
 \zm(t)P^*_{q,t}-\zm(t)E^*_{q,t}+J^*_{q,t}
 =\sum_{u\in I_x}F(u)\zm(u).
\]
The operator deletes $u=t$, where
$v_{q,t}(t)=(q-2)/(q-1)$.  Multiplication by the outer $q$ and summation prove
\eqref{eq:full}.

For the Poisson attachment, use the reflected-conjugate profile
$\phi(v)=\overline{\psi_+(-v)}$.  Reflection and conjugation preserve support
in $[-1,1]$, and its physical kernel is $\overline{K_H(-h)}$.  Thus the V43
complete-lattice theorem \cite{v43} applies in the adjoint direction without
assuming that $K_H$ is even or real.  When $q\leq2Q<H$, every nonzero dual
frequency lies outside the profile support, so $P^*_{q,t}=0$.  Finally,
\[
 \zm|_L-\zm|_R=\rho\left(\frac1\ell+\frac1r\right)=\rho^{-1},
\]
and only the opposite child survives in \eqref{eq:j}, proving
\eqref{eq:jump}.
\end{proof}

\section{The literal beta pairing and the returned shell coefficient}

The counting inner product is conjugate-linear in its first slot.  Since
$\beta$ and $\zm$ are real,
\begin{equation}
 \ip{\zm}{\Ax\beta}=\ip{\Ax^*\zm}{\beta}.                       \label{eq:adjoint}
\end{equation}
Let $E_{q,t}$ and $J_{q,t}$ denote the conjugates of the starred quantities in
\eqref{eq:e}--\eqref{eq:j}.  Under $H>2Q$, Theorem~\ref{thm:compiler} yields
\begin{align}
 \ip{\zm}{\Ax\beta}={}&D_{\beta,\zm}
 -\sum_{q\in\Qx}q\sum_{\substack{t\in I_x\\q\nmid t}}
       \beta(t)\zm(t)E_{q,t} \notag\\
 &+\sum_{q\in\Qx}q\sum_{\substack{t\in I_x\\q\nmid t}}
       \beta(t)J_{q,t},                                         \label{eq:scalar}
\end{align}
where
\begin{equation}
 D_{\beta,\zm}=-\sum_{q\in\Qx}\frac{q(q-2)}{q-1}
 \sum_{\substack{t\in I_x\\q\nmid t}}\zm(t)\beta(t).         \label{eq:d}
\end{equation}
Put
\begin{equation}
 B_Q=\sum_{q\in\Qx}\frac{q(q-2)}{q-1}.
\end{equation}
Restoring the omitted input multiples gives the exact identity
\begin{equation}
 D_{\beta,\zm}=-B_Q\ip{\zm}{\beta}
 +\sum_{q\in\Qx}\frac{q(q-2)}{q-1}
   \sum_{\substack{t\in I_x\\q\mid t}}\zm(t)\beta(t).         \label{eq:bq}
\end{equation}
The coefficient is the same deleted-diagonal shell coefficient isolated in
V43.  Equations \eqref{eq:scalar}--\eqref{eq:bq} are an exact arithmetic
normal form, not an estimate.

\section{Unit-mask and normalization firewalls}

For $q\nmid t$, write
\begin{equation}
 c_{q,t}(u)=\one_{u\equiv t\,(q)}-\frac1{q-1},\qquad
 d_q(u)=\frac{\one_{q\mid u}}{q-1}.                              \label{eq:cd}
\end{equation}
Then $v_{q,t}=c_{q,t}+d_q$, but over one period
\begin{equation}
 \sum_{u\bmod q}c_{q,t}(u)=-\frac1{q-1},\qquad
 \sum_{u\bmod q}d_q(u)=+\frac1{q-1}.                            \label{eq:modes}
\end{equation}
Only their sum is centered.  Their complete-lattice zero-frequency terms are
$-H\psi_+(0)/(q(q-1))$ and $+H\psi_+(0)/(q(q-1))$; applying centered
Poisson after discarding either term is invalid.  This residue-class
normalization is distinct from the deleted-diagonal value
$K_H(0)=\widehat\psi_+(0)=\int\psi_+=1$.  If $q\mid t$, the input mask in
\eqref{eq:operator} makes the entire $q$ summand zero, and no centered row is
constructed.

\section{Finite validation and adversarial controls}

The executable certificate is deliberately finite and exact.  Its primary
fixture uses the literal clock $x=64$, so $I_x=\{33,\ldots,64\}$,
$\ell=r=16$, $\rho^2=8$, and $\Qx=\{5,7\}$.  To avoid introducing an
irrational representation of $\rho$, the executable checks the rational step
$h=\zm/\rho$; linearity then gives the identical formula for $\zm$.  It uses
the literal $\beta$, together with a nonreal, noneven Gaussian-rational kernel
supported on the nineteen differences $-9,\ldots,9$ and normalized by
$K(0)=1$.  The producer computes $\Ax^*h$ directly and reassembles all 32
coordinates from the four terms in \eqref{eq:full}.  It also checks
\[
 \ip h{\Ax\beta}=\ip{\Ax^*h}{\beta}
\]
and the two equivalent forms of the deleted-diagonal lane.  Thus a missing
conjugate, outer $q$, mask, diagonal sign, or $q-2$ coefficient is visible in
exact arithmetic.

The independent checker imports no producer code.  It reconstructs the
certificate independently, verifies eight frozen source hashes and strict
integer types at nested fields, and rejects 100 adversarial mutations in 14
named classes.  The separate stress program checks 192 deterministic exact
families, split evenly between integer and noninteger clocks.  Its coverage
ledger contains 8448 coordinate decompositions, 9690 left-child and 9944
right-child jump checks, 5710 input-mask rows, and 206110 output-mask terms;
all 192 kernels are nonreal and noneven.  Normal and optimized Python
executions have byte-identical output and empty standard error.

These computations do not approximate $\psi_+$, take a large-$x$ limit, or
test prime cancellation.  The complete-row zero in Theorem~\ref{thm:compiler}
is supported by the frozen Poisson theorem, not inferred from the finite
fixture.  Conversely, the fixture retains a generally nonzero $P^*_{q,t}$ in
\eqref{eq:full}, so the algebraic decomposition is tested before the
source-backed zero is substituted.  A separate complete-period row checks
\eqref{eq:modes} and rejects dropping the output-unit correction.

\section{Route status}

TPC-255 proves exact literal structure.  It supplies no estimate for
$\ip{\zm}{\beta}$, the input-unit correction, or the two boundary lanes.
Triangulating these pieces prime by prime would remove precisely the signed
reassembly one must eventually exploit.  Accordingly
\[
 \texttt{ROUTE\_ADVANCE=YES\_EXACT\_LITERAL\_STRUCTURE},
\]
\[
 \texttt{LITERAL\_ARITHMETIC\_STRUCTURE\_ADVANCE=YES},\qquad
 \texttt{ARITHMETIC\_ADVANCE=NO}.
\]
The strict $1/400$ budget is unpaid and no twin-prime statement follows.
The next admissible theorem must estimate the same signed object in
\eqref{eq:scalar}--\eqref{eq:bq}, with both masks and all three surviving
lanes retained.

\paragraph{Strongest result and obstruction.}
The complete-lattice centered alias has been removed from the literal
adjoint Haar form without an evenness or self-adjointness assumption.  What
remains is an exact three-lane object: the $B_Q$-weighted local return, the
outer hard-window leakage, and the internal child-jump leakage, together with
the input-unit correction already contained in the first lane.  This is the
strongest positive conclusion.  The matching obstruction is equally exact:
Poisson cancellation acts before diagonal deletion and therefore cannot pay
the returned $B_Q$ term.  No frozen theorem controls that term collectively
with both boundary lanes on the V59 clock.

\paragraph{Reusable compiler and open theorem.}
The reusable sequence is
\[
 \text{literal adjoint test}
 \longrightarrow \text{unit-centered complete lattice}
 \longrightarrow \text{Poisson zero}
 \longrightarrow \text{diagonal and boundary return}.
\]
The open theorem is a signed estimate for \eqref{eq:scalar} after the
substitution \eqref{eq:bq}, preserving the whole prime shell, both unit masks,
and both physical boundaries.  Separate absolute estimates for the displayed
pieces would not establish that theorem.  Thus TPC-255 identifies the next
arithmetic target exactly while leaving $L^2$, fixed-atom credit, full Gate B,
and the twin-prime problem open.

\bibliographystyle{plain}
\bibliography{references}
\end{document}
